\section{Рабочий проект}

\subsection{Маршруты и компоненты программной системы}

Рабочий проект платформы Easyappoint представлен тремя основными программными компонентами: клиентским web-приложением, серверным API-сервисом и сервисом Telegram-уведомлений. Компоненты разрабатываются независимо, но используют общие контракты REST-ответов и Kafka-событий.

Клиентское приложение находится в каталоге \texttt{apps/frontend}. Оно реализовано на React и TypeScript и содержит маршрутизацию для публичной страницы, кабинета мастера, кабинета клиента и административного раздела. Серверная часть находится в каталоге \texttt{apps/data-service}. Она реализована на Kotlin с использованием Ktor и отвечает за аутентификацию, пользователей, расписание, приглашения, бронирования, интеграцию с календарем и публикацию событий. Сервис уведомлений находится в каталоге \texttt{apps/telegram-bot}; он читает события из Kafka и отправляет пользователям Telegram-сообщения.

Основные маршруты клиентского приложения представлены в таблице \ref{frontend:routes}.

\begin{xltabular}{\textwidth}{|p{4.2cm}|p{4.6cm}|X|}
\caption{Основные маршруты клиентского приложения\label{frontend:routes}}\\ \hline
\centrow Маршрут & \centrow Компонент & \centrow Назначение \\ \hline
\endfirsthead
\caption*{Продолжение таблицы \ref{frontend:routes}}\\ \hline
\centrow Маршрут & \centrow Компонент & \centrow Назначение \\ \hline
\finishhead
\texttt{/} & \texttt{PublicLandingPage} & Публичная страница входа в платформу. \\ \hline
\texttt{/login}, \texttt{/register} & \texttt{LoginPage}, \texttt{RegisterPage} & Аутентификация и регистрация пользователя. \\ \hline
\texttt{/link} & \texttt{LinkLandingPage} & Переход клиента по пригласительной ссылке мастера. \\ \hline
\texttt{/master/appointments} & \texttt{MasterAvailabilityPage} & Основной экран мастера для управления расписанием и записями. \\ \hline
\texttt{/master/clients} & \texttt{MasterClientsPage} & Работа мастера с клиентами и приглашениями. \\ \hline
\texttt{/client/availability} & \texttt{ClientAvailabilityPage} & Просмотр клиентом доступных слотов мастера. \\ \hline
\texttt{/client/appointments} & \texttt{ClientBookingsPage} & Просмотр и отмена записей клиента. \\ \hline
\texttt{/admin} & \texttt{AdminOverviewPage} & Сводная административная панель состояния платформы. \\ \hline
\texttt{/admin/users} & \texttt{AdminUsersPage} & Просмотр пользователей, фильтрация и изменение их статуса и ролей. \\ \hline
\texttt{/admin/appointments} & \texttt{AdminAppointmentsPage} & Просмотр и административное управление записями. \\ \hline
\texttt{/admin/audit} & \texttt{AdminAuditPage} & Просмотр журнала административных действий. \\ \hline
\texttt{/admin/monitoring} & \texttt{AdminMonitoringPage} & Просмотр эксплуатационного состояния системы администратором. \\ \hline
\end{xltabular}

Основные серверные маршруты, используемые в рабочих сценариях, представлены в таблице \ref{backend:routes}.

\begin{xltabular}{\textwidth}{|p{4.2cm}|p{4.3cm}|X|}
\caption{Основные маршруты серверного API\label{backend:routes}}\\ \hline
\centrow Маршрут & \centrow Контроллер & \centrow Назначение \\ \hline
\endfirsthead
\caption*{Продолжение таблицы \ref{backend:routes}}\\ \hline
\centrow Маршрут & \centrow Контроллер & \centrow Назначение \\ \hline
\finishhead
\texttt{/auth/*} & \texttt{PasswordAuthController} & Регистрация, вход в систему, обновление и завершение сессии. \\ \hline
\texttt{/users/me} & \texttt{UserMeController} & Получение данных текущего пользователя и его ролей. \\ \hline
\texttt{/availability/*} & \texttt{AvailabilityController} & Получение доступных слотов и настройка расписания мастера. \\ \hline
\texttt{/invites/*} & \texttt{InviteController} & Создание, получение и изменение статуса пригласительных ссылок. \\ \hline
\texttt{/appointments/*}, \texttt{/bookings/*} & \texttt{BookingController} & Создание, изменение, отмена и просмотр бронирований. \\ \hline
\texttt{/google/*} & контроллер Google Calendar & Подключение внешнего календаря и синхронизация занятости. \\ \hline
\texttt{/notifications/*} & \texttt{NotificationController} & Управление каналами уведомлений пользователя. \\ \hline
\texttt{/admin/*} & \texttt{AdminController} & Административный обзор, управление пользователями и записями, аудит и мониторинг. \\ \hline
\end{xltabular}

\subsection{Спецификация классов программной системы}

Ключевые классы и модули программной системы представлены в таблице \ref{class:table}. Они выделены по фактической структуре исходного кода и отражают основные рабочие сценарии платформы.

\renewcommand{\arraystretch}{0.85}
\begin{xltabular}{\textwidth}{|p{5cm}|p{3.6cm}|X|}
\caption{Описание основных классов и модулей Easyappoint\label{class:table}}\\ \hline
\centrow Класс или модуль & \centrow Расположение & \centrow Назначение \\ \hline
\endfirsthead
\caption*{Продолжение таблицы \ref{class:table}}\\ \hline
\centrow Класс или модуль & \centrow Расположение & \centrow Назначение \\ \hline
\finishhead
\texttt{BookingController} & \texttt{apps/data-service} & Описывает REST-маршруты для создания, изменения, отмены и просмотра бронирований. \\ \hline
\texttt{BookingService} & \texttt{apps/data-service} & Содержит бизнес-логику бронирований: проверку роли пользователя, проверку пересечений, создание связи мастер-клиент, запись событий в outbox и синхронизацию с календарем. \\ \hline
\texttt{AvailabilityService} & \texttt{apps/data-service} & Рассчитывает доступность мастера с учетом правил расписания, активных бронирований и занятости во внешнем календаре. \\ \hline
\texttt{AvailabilitySlotCalculator} & \texttt{apps/data-service} & Строит кандидатные временные слоты, разбивает интервалы по длительности записи и исключает занятые окна. \\ \hline
\texttt{InviteService} & \texttt{apps/data-service} & Управляет жизненным циклом пригласительных ссылок и проверяет доступ клиента к мастеру. \\ \hline
\texttt{GoogleCalendarSyncService} & \texttt{apps/data-service} & Выполняет синхронизацию с Google Calendar и поддерживает данные о внешней занятости мастера. \\ \hline
\texttt{BookingOutboxWriter} & \texttt{apps/data-service} & Формирует события \texttt{booking.created}, \texttt{booking.updated}, \texttt{booking.cancelled} для дальнейшей публикации в Kafka. \\ \hline
\texttt{TelegramBotEventProcessor} & \texttt{apps/telegram-bot} & Обрабатывает Kafka-события и вызывает сервисы отправки Telegram-уведомлений. \\ \hline
\texttt{AdminController} & \texttt{apps/data-service} & Описывает защищенные REST-маршруты администратора для пользователей, записей, аудита и мониторинга. \\ \hline
\texttt{AdminUserService} & \texttt{apps/data-service} & Реализует административный обзор пользователей, фильтрацию и изменение статуса и ролей. \\ \hline
\texttt{AdminAppointmentService} & \texttt{apps/data-service} & Реализует просмотр, изменение и отмену записей от имени администратора. \\ \hline
\texttt{App} & \texttt{apps/frontend} & Определяет маршруты React-приложения и ролевую навигацию для мастера, клиента и администратора. \\ \hline
\texttt{ClientAvailabilityPage} & \texttt{apps/frontend} & Реализует экран выбора доступного слота клиентом. \\ \hline
\texttt{MasterAvailabilityPage} & \texttt{apps/frontend} & Реализует рабочий экран мастера для управления доступностью и записями. \\ \hline
\end{xltabular}
\renewcommand{\arraystretch}{1.0}

Рабочая логика создания бронирования сосредоточена в классе \texttt{BookingService}. При создании записи сервис проверяет роль пользователя, корректность временного интервала, существование мастера, допустимость пригласительной ссылки, отсутствие пересечений с активными записями и доступность выбранного слота. После успешного сохранения создается связь между мастером и клиентом, записывается outbox-событие и запускается синхронизация с Google Calendar.

Расчет доступных слотов выделен в отдельный класс \texttt{AvailabilitySlotCalculator}. Такое разделение упрощает тестирование алгоритма: кандидатные слоты строятся из регулярного расписания мастера и импортированных окон доступности, после чего из них исключаются занятые интервалы.

Для обмена между серверным API и сервисом уведомлений используется outbox-подход, соответствующий распространенному шаблону transactional outbox \cite{transactionalOutbox}. Серверная часть сохраняет доменное событие в таблицу \texttt{outbox\_events}, после чего отдельный publisher публикует его в Kafka. Сервис \texttt{telegram-bot} получает событие, восстанавливает payload и отправляет уведомление пользователю.

\subsection{Реализация основных пользовательских сценариев}

Рабочий сценарий платформы начинается с регистрации пользователя. На серверной стороне запросы регистрации и входа обрабатываются сервисом парольной аутентификации: пароль сохраняется только в виде хеша, а успешная авторизация возвращает JWT-токены. При каждом обращении к защищенному маршруту система восстанавливает текущего пользователя и его роли, причем фактические роли берутся из базы данных, а не только из содержимого токена.

Сценарий мастера реализован вокруг настройки рабочей доступности. Мастер задает часовой пояс, длительность приема и набор правил расписания по дням недели. При сохранении расписания сервер проверяет пересечения правил одного дня, обновляет настройки мастера и синхронизирует зеркальные интервалы с внешним календарем при активном подключении Google Calendar.

Сценарий клиента строится через приглашение. Мастер создает пригласительную ссылку, после чего клиент переходит по ней, регистрируется или входит в систему и получает связь с мастером. При выборе времени клиентский интерфейс запрашивает доступные слоты, а сервер рассчитывает их на основании расписания мастера, активных бронирований и занятости из Google Calendar.

Создание бронирования выполняется как атомарная бизнес-операция. Сервис проверяет допустимость временного окна, связь клиента с мастером, отсутствие пересечения с другими активными записями и доступность выбранного слота. После сохранения записи создается событие для уведомлений, а при подключенном календаре выполняется синхронизация внешнего события. Отмена и перенос записи проходят через те же проверки доступа и также создают события для последующей доставки уведомлений.

Административный сценарий вынесен в отдельный защищенный раздел. Администратор получает сводку по пользователям и записям, просматривает аудит действий, изменяет статусы пользователей и управляет записями при необходимости ручного вмешательства. Критичные административные операции фиксируются в таблице аудита с указанием пользователя-исполнителя, целевого объекта и контекста изменения.

\subsection{Тестирование разработанной платформы}

Тестирование платформы выполняется на уровне серверной бизнес-логики, клиентского интерфейса, API-оберток, контрактов и интеграционных сценариев. Поддерживаемая проверка репозитория запускается командой \texttt{make ci}; она выполняет проверку контрактов, frontend lint/test/build, backend Gradle check, telegram-bot Gradle check и render Kubernetes-манифестов.

Структура проверок представлена в таблице \ref{testing:levels}.

\begin{xltabular}{\textwidth}{|L{3.4cm}|L{4.2cm}|Y|}
\caption{Уровни тестирования платформы Easyappoint\label{testing:levels}}\\ \hline
\centrow Уровень проверки & \centrow Используемые средства & \centrow Назначение \\ \hline
\endfirsthead
\caption*{Продолжение таблицы \ref{testing:levels}}\\ \hline
\centrow Уровень проверки & \centrow Используемые средства & \centrow Назначение \\ \hline
\finishhead
Серверные unit-тесты & JUnit, kotlin.test, Mockito & Проверка правил записей, расписания, приглашений, JWT-аутентификации, календаря и outbox-публикации. \\ \hline
Frontend-тесты & Vitest, React Testing Library & Проверка страниц входа, кабинетов, календарей, форм и клиентских API-сценариев. \\ \hline
Контрактные тесты & JSON Schema Validator, \texttt{packages/contracts} & Проверка совместимости REST DTO и Kafka-событий с опубликованными контрактами \cite{jsonSchema}. \\ \hline
Интеграционные сценарии & E2E-поток и сервисные интеграционные тесты & Проверка приглашения, записи, отмены, генерации outbox-событий и обработки уведомлений. \\ \hline
Инфра-проверка & \texttt{make ci}, Gradle, npm, kustomize render & Проверка сборки компонентов и render Kubernetes-манифестов перед развертыванием. \\ \hline
\end{xltabular}

Основные тестовые сценарии представлены в таблице \ref{testing:table}.

\begin{xltabular}{\textwidth}{|L{3.6cm}|L{3.9cm}|Y|}
\caption{Основные сценарии тестирования Easyappoint\label{testing:table}}\\ \hline
\centrow Сценарий & \centrow Проверяемый модуль & \centrow Ожидаемый результат \\ \hline
\endfirsthead
\caption*{Продолжение таблицы \ref{testing:table}}\\ \hline
\centrow Сценарий & \centrow Проверяемый модуль & \centrow Ожидаемый результат \\ \hline
\finishhead
Регистрация мастера и вход & Тесты аутентификации и страницы входа & Пользователь регистрируется с ролью мастера; пароль хранится хешем, вход возвращает сессию. \\ \hline
Создание записи с пересечением & Тесты сервиса бронирований & Система возвращает конфликт и не сохраняет некорректную запись. \\ \hline
Создание записи вне доступности & Тесты сервиса бронирований & Система отклоняет запись, если слот отсутствует в расписании. \\ \hline
Создание записи без приглашения & Тесты сервиса бронирований & Запись разрешена только клиенту, связанному с мастером. \\ \hline
Перенос активной записи & Тесты бронирований и страницы мастера & Система обновляет окно, проверяет пересечения и публикует событие изменения. \\ \hline
Расчет доступных слотов & Тесты доступности и интеграционного сценария & Слоты строятся из расписания, затем исключаются активные записи и занятость Google Calendar. \\ \hline
Настройка расписания мастера & Тесты расписания и страницы мастера & Система сохраняет правила недели, отклоняет пересечения и показывает слоты. \\ \hline
Жизненный цикл приглашения и записи & Frontend e2e API & Клиент получает приглашение, создает запись и отменяет ее через ожидаемые API-вызовы. \\ \hline
Проверка REST-контрактов & Контрактные тесты серверного API & Фактические DTO соответствуют JSON-схемам из \texttt{packages/contracts}. \\ \hline
Проверка Kafka-контрактов & Контрактные тесты событий & События уведомлений соответствуют опубликованным схемам. \\ \hline
Синхронизация Google Calendar & Тесты календарной интеграции & Создание, отмена записи и зеркала расписания преобразуются во внешние операции. \\ \hline
Подключение Telegram & Тесты сервиса уведомлений и bot-link сценария & Пользователь получает ссылку, подтверждает ее в боте, а система сохраняет назначение. \\ \hline
Обработка Telegram-события & Тесты обработчика telegram-bot & Сервис уведомлений определяет тип события и вызывает нужный обработчик. \\ \hline
Повтор Kafka-события & Тесты record handler telegram-bot & При временной ошибке выполняются повторы, после лимита запись уходит в DLT. \\ \hline
Аудит администратора & Тесты административного сервиса и страницы аудита & Действия доступны для просмотра с фильтрами по пользователю, действию и объекту. \\ \hline
\end{xltabular}

Результаты проверок по ключевым рабочим сценариям представлены в таблице \ref{testing:results}.

\begin{xltabular}{\textwidth}{|L{3.5cm}|L{4.3cm}|Y|}
\caption{Результаты проверки ключевых сценариев\label{testing:results}}\\ \hline
\centrow Проверяемый сценарий & \centrow Критерий успешности & \centrow Результат \\ \hline
\endfirsthead
\caption*{Продолжение таблицы \ref{testing:results}}\\ \hline
\centrow Проверяемый сценарий & \centrow Критерий успешности & \centrow Результат \\ \hline
\finishhead
Регистрация и аутентификация & Создание пользователя, выдача токена, восстановление профиля и ролей & Реализовано и покрыто серверными и frontend-тестами. \\ \hline
Настройка расписания мастера & Сохранение правил недели, контроль пересечений, учет часового пояса & Реализовано, проверяется сервисом и страницей мастера. \\ \hline
Выбор свободного слота & Отображение доступности с учетом расписания, записей и внешнего календаря & Реализовано, проверяется тестами доступности и клиентского календаря. \\ \hline
Создание и отмена записи & Проверка прав, сохранение записи, смена статуса, публикация событий & Реализовано, проверяется unit- и e2e-сценариями. \\ \hline
Telegram-уведомления & Подключение канала, обработка Kafka-события, отправка сообщения & Реализовано, проверяется сервисом уведомлений и telegram-bot. \\ \hline
Google Calendar & Создание, удаление и синхронизация внешних календарных событий & Реализовано, проверяется клиентом и сервисом синхронизации. \\ \hline
Админ-контроль & Просмотр пользователей, записей, аудита и эксплуатационного состояния & Реализовано в отдельном административном разделе. \\ \hline
\end{xltabular}

В ходе проверки были охвачены основные пользовательские сценарии платформы: регистрация и аутентификация пользователя, настройка расписания мастера, создание пригласительной ссылки, выбор клиентом свободного слота, создание и отмена бронирования, отправка уведомлений, синхронизация с Google Calendar и административный контроль.

Проверка показала, что разработанная система корректно объединяет клиентский интерфейс, серверную бизнес-логику, хранение данных, внешние интеграции и автоматическую проверку основных сценариев. Таким образом, рабочий проект соответствует функциональным требованиям, сформулированным в техническом задании.
